\chapter{实验设计与正确性验证}\label{chap:verification}

\section{实验目标与总体设计}

第 3 章给出了 LoongArch 平台 ART 执行引擎的适配设计与实现，第 5 章则将进一步讨论三项正式优化
方案的性能结果。本章的任务，是说明这些结论由怎样的实验设计支撑，并回答两个问题：其一，适配后的
执行引擎是否具备可靠的语义正确性；其二，后续性能结果是否建立在一致、可重复的实验条件之上。

围绕上述目标，本文将实验分为正确性验证和性能评测两部分。所有实验统一建立在同一 LoongArch64
目标设备和同一源码状态下，并遵循如下顺序：
\begin{enumerate}
    \item 完整编译当前 ART 产物；
    \item 采用轻量同步方式部署到目标机；
    \item 执行目标机回归与类库验证；
    \item 在验证通过后运行基准测试，采集分数、日志与热点数据。
\end{enumerate}

\begin{figure}[htbp]
    \centering
    \includegraphics[width=0.95\textwidth]{verification/experiment-validation-workflow.pdf}
    \caption{本文正确性验证与性能评测的组织流程}
    \label{fig:verification-workflow}
\end{figure}

图~\ref{fig:verification-workflow} 概括了从固定源码状态到正式数据归档的实验组织流程，其中“先正确性，
后性能”是全文实验设计的基本约束。这一组织方式的核心目的，是将“系统正确”和“性能变化”严格区分。若相关回归未通过，则后续跑分结果不具有
解释价值；只有在正确性得到约束后，性能差异才可以归因于方案本身。本章因此不再展开与本地环境强绑定的
脚本细节，而重点说明支撑论文结论的实验思路、数据来源和验证结果。

\section{实验环境与执行设置}

本文的实验平台基于真实 LoongArch64 目标设备，处理器为龙芯 3A6000，四核，主频
\texttt{2.5\,GHz}，内存为 \texttt{16\,GB}；目标系统为 Android 15，内核版本为
Linux 6.6.46。软件侧以 AOSP 15 源码中的 \texttt{art/} 子项目为研究对象，在完成第 3 章所述
平台适配后开展后续验证与性能评测。

本文的正式性能结论统一建立在 LoongArch64 64 位 \texttt{debug} ART 与 \texttt{JIT} 模式下。
选择这一设置有两方面原因。第一，\texttt{debug} 构建更容易暴露调用约定、异常恢复和慢路径跳转等
平台相关错误，有利于在适配尚未完全稳定时完成收敛。第二，第 5 章的三项正式方案也都在相同设置下完成
对比，因此能够保证前后章节的实验条件一致。

ART 改动统一通过 \texttt{art/tools/buildbot-build.sh --target} 完整编译，之后采用轻量同步方式
将相关产物部署到目标机。保留这一执行链路的目的，在于保证待测产物与当前源码状态一致，而不是将论文
正文写成操作手册式的环境说明。

\section{正确性验证设计与结果}

\subsection{分层验证思路}

本文采用分层验证方式约束正确性，具体包括以下三个层次。
\begin{enumerate}
    \item \textbf{宿主机侧测试。} 用于尽早发现与平台化改动相关的公共组件回归，例如 DEX 处理、
    命令行参数解析和编译器基础设施错误。
    \item \textbf{目标机侧 run-test。} 用于验证解释执行、编译执行、异常处理、对象访问和调用约定等
    执行引擎核心语义。
    \item \textbf{libcore 测试。} 用于进一步验证 Java 核心类库层面的 API 行为是否正常。
\end{enumerate}

这种分层结构的意义在于：宿主机侧测试负责前置过滤，目标机侧 run-test 负责运行时语义验收，
libcore 测试则提供更接近 Java 应用视角的补充证据。前一层通过后，才进入下一层验证与后续性能评测。

\subsection{目标机侧 run-test 整体验收}

目标机侧 run-test 是本文验证执行引擎正确性的核心手段。整体验收使用如下命令：
\begin{lstlisting}[language=bash,basicstyle=\ttfamily\small]
art/test.py -r --target --no-prebuild --debug --no-image --64 --all-compiler -j 7
\end{lstlisting}

该命令固定了目标机、\texttt{debug}、\texttt{no-prebuild}、\texttt{no-image} 和 64 位等
关键条件，并通过 \texttt{--all-compiler} 展开 \texttt{interp-ac}、\texttt{interpreter}、
\texttt{jit}、\texttt{jit-on-first-use}、\texttt{optimizing}、\texttt{speed-profile}
和 \texttt{baseline} 七种 compiler 变体。由于 \texttt{art/test} 目录当前包含
\texttt{1033} 个 run-test 用例，因此理论配置总数为 \texttt{7231}。测试框架会依据
\texttt{knownfailures.json} 中的规则自动跳过不适用配置，最终执行结果如
表~\ref{tab:target-run-test-result} 所示。

\begin{table}[htbp]
    \centering
    \caption{目标机侧 run-test 验证结果}
    \label{tab:target-run-test-result}
    \begin{tabular}{lccc}
        \hline
        \textbf{compiler 变体} & \textbf{理论配置数} & \textbf{自动 skip 数} & \textbf{实际执行数} \\
        \hline
        \texttt{interp-ac} & 1033 & 67 & 966 \\
        \texttt{interpreter} & 1033 & 59 & 974 \\
        \texttt{jit} & 1033 & 57 & 976 \\
        \texttt{jit-on-first-use} & 1033 & 61 & 972 \\
        \texttt{optimizing} & 1033 & 70 & 963 \\
        \texttt{speed-profile} & 1033 & 59 & 974 \\
        \texttt{baseline} & 1033 & 57 & 976 \\
        \hline
        \textbf{总计} & \textbf{7231} & \textbf{430} & \textbf{6801} \\
        \hline
    \end{tabular}
\end{table}

表~\ref{tab:target-run-test-result} 表明，LoongArch 平台 ART 已能够在固定目标机条件下对
\texttt{1033} 个 run-test 用例和 \texttt{7} 类 compiler 变体进行系统级回归。其中，
七种 compiler 变体是 ART run-test 的测试组织维度，并不与前文按实现机制划分的 C++ 解释器、
Nterp 和 JIT 三条主路径一一对应。整体结果说明，第 3 章建立的执行引擎能力已经具备进入正式性能
评测阶段的语义基础。

\subsection{libcore 测试结果}

仅依赖 run-test 仍不足以证明 Java 核心类库层面的行为完全正确，因此本文进一步引入 libcore
测试集，用于覆盖集合框架、输入输出、并发、数学运算、加解密与网络等 Java API 语义。按当前实验
归档保留的包级摘要，设备侧已覆盖 libcore 的主要测试包，验证结果以通过为主，少量测试因架构条件
限制被自动跳过。现有记录保留的是包级摘要而非逐包统计明细，因此本文不再将其写成精确通过率数字；
这些自动跳过项不涉及执行引擎核心语义，不影响主要正确性结论。

\subsection{优化阶段的最小相关回归}

性能优化阶段并不重复执行全量系统回归，而是遵循“最小相关回归优先”的原则。即每个候选优化方案在进入
基准测试之前，先根据改动所在模块选择直接受影响的目标机测试集合。本文三项正式方案采用了如下代表性
验证方式：
\begin{itemize}
    \item \textbf{方案一（JIT 分层热度阈值调整）。} 重点验证 \texttt{708-jit-cache-churn}、
    \texttt{667-jit-jni-stub} 等 JIT 相关用例。
    \item \textbf{方案二（运行时字符串搬运批量化）。} 重点验证
    \texttt{624-checker-stringops}、\texttt{729-checker-polymorphic-intrinsic}、
    \texttt{595-profile-saving}、\texttt{570-checker-osr}、
    \texttt{667-jit-jni-stub} 和 \texttt{708-jit-cache-churn}。
    \item \textbf{方案三（LoongArch64 LSX 浮点 SIMD 向量化）。} 重点验证
    \texttt{527-checker-array-access-simd}、\texttt{656-checker-simd-opt}、
    \texttt{570-checker-osr} 以及 \texttt{708-jit-cache-churn}。
\end{itemize}

这种做法的目的，是在保证验证效率的同时维持结论可信度。若相关回归未通过，则基准测试分数不具有
解释价值；只有在最小相关回归通过后，后续性能差异才可以归因于方案本身，而不是功能错误或未定义行为。

\section{性能评测设计}

\subsection{基准测试与工作负载选择}

性能评测的目标，不是给出可直接外推到所有 Android 应用场景的绝对结论，而是在统一实验条件下观察
执行引擎优化对代表性工作负载的影响。基于这一目标，本文选择 DaCapo 与 SPECjvm2008 作为主要
基准测试集合\citep{Blackburn2006,SPECjvm2008}。前者更偏向对象分配、字符串处理与复杂控制流；
后者则包含更接近经典 JVM 数值与吞吐路径的工作负载。已有研究也表明，不同基准测试子项在分配速率、
存活集规模和 GC 行为上存在显著差异，因此基准选择本身就是实验设计的一部分
\citep{Lengauer2017GCBench,Prokopec2019Renaissance}。

结合热点分析结果，本文在正式图表中保留以下五个工作负载：
\begin{itemize}
    \item \texttt{DaCapo/pmd}：用于观察 JIT 进入时机、解释器驻留以及字符串运行时热点；
    \item \texttt{SPECjvm2008/mpegaudio}：用于提供与 \texttt{pmd} 明显不同的吞吐型对照；
    \item \texttt{SPECjvm2008/scimark.lu.small}：用于验证 LoongArch64 LSX 浮点 SIMD 方案；
    \item \texttt{SPECjvm2008/scimark.sor.small}：用于观察数值型优化的可迁移性；
    \item \texttt{SPECjvm2008/scimark.fft.small}：用于补充复杂数值内核的边界对照。
\end{itemize}

其中，DaCapo 的评价指标为执行耗时，单位为 \texttt{msec}，数值越小表示性能越好；SPECjvm2008
的评价指标为 \texttt{ops/m}，数值越大表示吞吐越高。本文在后续分析中始终遵循这一指标方向，不将
不同基准测试的分值含义混为一谈。

\subsection{正式对比数据集与重复测量数据集}

为了控制设备状态和日期漂移带来的干扰，本文将正式单次对比与重复测量验证明确分离。

\begin{itemize}
    \item \textbf{正式对比数据集：\texttt{test-data/final-comparison/}}。该数据集采集于
    2026 年 3 月 13 日，用于正文主表和主图，包含同日 baseline、三项正式方案以及叠加状态。
    本文后续单次正式结果若无特别说明，均指这一数据集。其基线结果如
    表~\ref{tab:formal-baseline-values} 所示。
    \item \textbf{重复测量数据集：\texttt{test-data/statistical-validation-20260314/}}。
    该数据集采集于 2026 年 3 月 14 日，仅用于验证两条最关键的性能主张，即方案一在
    \texttt{pmd} 上的收益方向和方案三在 \texttt{lu.small} 上的稳定收益。其结果单独展示，
    不与前一日的正式对比数据混算。
    \item \textbf{早期 baseline 归档：\texttt{test-data/final-baseline/}}。该目录保留的是较早
    阶段的对照结果，用于历史记录，不再作为三项正式方案总对比的主数据源。
\end{itemize}

\begin{table}[htbp]
    \centering
    \small
    \caption{正式对比数据集的五个基线结果}
    \label{tab:formal-baseline-values}
    \begin{tabular}{lc}
        \hline
        \textbf{工作负载} & \textbf{基线结果} \\
        \hline
        \texttt{DaCapo/pmd} & \texttt{3272\,msec} \\
        \texttt{SPECjvm2008/mpegaudio} & \texttt{25.38\,ops/m} \\
        \texttt{SPECjvm2008/scimark.lu.small} & \texttt{46.22\,ops/m} \\
        \texttt{SPECjvm2008/scimark.sor.small} & \texttt{24.76\,ops/m} \\
        \texttt{SPECjvm2008/scimark.fft.small} & \texttt{76.43\,ops/m} \\
        \hline
    \end{tabular}
\end{table}

通过将单次正式对比与重复测量验证分离，本文能够避免把跨日期、跨设备状态的变化误写成方案收益，
并使第 5 章对“单次结果”和“稳定结论”的区分具有明确的数据基础。

\subsection{数据采集与分析方法}

为保证性能分析不仅停留在单一分数层面，本文在基准测试运行时同步保留如下原始产物：
\begin{itemize}
    \item \texttt{run.log}：记录脚本调用、远端命令和运行参数；
    \item \texttt{perf.data}：采样得到的原始性能数据；
    \item \texttt{report\_symbol.txt}：符号级热点汇总；
    \item \texttt{report\_callgraph.txt}：调用图视角的热点信息；
    \item \texttt{flamegraph.svg}：火焰图，用于直观展示热点分布。
\end{itemize}

这些数据使第 5 章可以同时从分数变化、热点分布变化和固定成本变化三个层面解释优化结果，而不是仅凭
单次跑分做出过度推断。正式图表所使用的数据、图文件和作图脚本，也统一保存在外部实验归档目录中，
便于后续交叉核验。

\section{有效性威胁与实验控制}

尽管本文已经通过统一设备、固定编译与部署顺序、最小相关回归和固定数据集组织方式控制实验变量，
面向 JIT 运行时的性能研究仍然不可避免地受到若干因素影响。为使第 5 章的结论不超出证据范围，
本文对以下因素予以明确约束。

\begin{enumerate}
    \item \textbf{设备状态与后台负载扰动。} 基准测试结果会受到目标机温度、后台服务活动、
    CPU 频率波动和前序运行历史的影响。因此，本文不混用不同日期的数据源，而是将
    \texttt{final-comparison} 与 \texttt{statistical-validation-20260314} 严格分离。
    \item \textbf{基准测试覆盖范围有限。} 本文正文最终保留的正式工作负载包括
    \texttt{pmd}、\texttt{mpegaudio}、\texttt{lu.small}、\texttt{sor.small} 和
    \texttt{fft.small}。这些工作负载已经能够支撑第 5 章的主要结论，但对并发类工作负载、
    复杂服务型应用和更长生命周期应用的覆盖仍然有限。因此，正文结果应被理解为固定实验设置下的
    阶段性结论，而不是 LoongArch 平台 ART 的完整性能画像。
    \item \textbf{\texttt{debug} 构建限制。} 正式性能结论统一建立在 LoongArch64
    \texttt{debug} ART 环境下。这一设置有利于适配阶段的问题暴露，但也会放大部分 JIT 调试信息与
    运行时固定成本。因此，本文所有性能结论都被限制为固定 LoongArch64 \texttt{debug JIT}
    条件下的阶段性结果，不直接外推到 \texttt{release} 构建或其它系统配置。
    \item \textbf{重复测量样本数有限。} 第 5 章仅对最关键的两条性能主张进行了 3 轮同日重复测量，
    其目的在于区分“可稳定复现的收益”和“仍需进一步验证的收益方向”，而不是进行严格的统计显著性检验。
    因此，论文使用原始散点、均值和标准差表达结论强度，而不滥用显著性术语。
    \item \textbf{JIT 符号可观测性有限。} 部分实验归档中存在 \texttt{perf-*.map} 缺失或
    JIT Java 符号不完整的情况。这意味着火焰图和 \texttt{report\_symbol.txt} 更适合用于判断
    “热点路径是否被命中”和“固定成本是否相对变化”，而不适合支撑过于细粒度的方法级归因。
\end{enumerate}

\section{本章小结}

本章从实验目标、验证层次、正式对比数据集和有效性威胁控制等方面，给出了本文的实验设计与正确性验证
方法。其核心作用不在于展开作者本地脚本细节，而在于说明第 3 章的适配结果如何得到语义正确性的支撑，
以及第 5 章的性能结论为何能够建立在一致、可重复的实验条件之上。
%---------------------------------------------------------------------------%
